% Transcription — fonds Alexandre Grothendieck, shelfmark 72, batch 6 (pages 101-112).
% Established from the facsimile archives/batches/72/batch-06.pdf (cut from a
% copy of 72.pdf fetched through the project's relay,
% https://grothendieck-relay.onrender.com/source/72.pdf, the scan published at
% https://grothendieck.umontpellier.fr/72.pdf), per the skill
% transcribe-grothendieck. The folder is 112 pages in six batches (1-20, 21-40,
% 41-60, 61-80, 81-100, 101-112), transcribed in parallel; this is the last.
% The raw scan's cover sheet is not in the batch file: PDF page k is
% archivists' page 100+k, and the pencilled numbers bottom-left (« 101 » ...
% « 112 ») agree on every page.
%
% Pass: Opus 5.5 (claude-opus-5-5), 2026-09-27 — first pass, unchecked against the pages by a human.
%
% Pages skipped: none. Page 106 is a tan wrapper sheet bearing only his
% heading, top right, « Classes de Stiefel-Whitney et classe de Brauer de la
% forme invariant »; it takes no \page marker and its words become the
% section heading, with a \note. Pages summarised: none. Pages 101-105 and
% 107-112 are his working notes, in blue-black ink, transcribed in full.
% No pagination of his own on these leaves; he numbers formulas and
% conditions instead, and the batch opens in the middle of such a run
% (references to (4), (6)-(9), (12), (13), (17)-(19), (24)-(26) point back to
% earlier batches; (27)-(29) are introduced here, (27) twice, once in a
% cancelled passage).
%
% What the batch is about. Not polytopes directly: two unrelated strands.
%   101-105  The end of a proof on chains of pseudo-reflections
%            z_i = id + a'_i (x) a_i in a locally free k-module E, with the
%            pairings <a_i, a'_j>: p. 101 closes (25), (25') with « cqfd »
%            (the relation z_i z_{i+1} - rho z_{i+1} z_i, a 2x2 triangular
%            matrix of determinant <a_{i+1},a'_i><a_i,a'_{i+1}> in k^x); p. 102
%            a Remarque for E of rank n+2 (conditions (A-bar), (B-bar), (26)
%            E_{n+1} = E) and a Proposition: for n+1 pseudo-reflections z_0..z_n
%            there is locally s_0 in E satisfying (A-bar), (B-bar) iff a) the
%            a_i are strictly free, a') so are the a'_i, b) <a_i,a'_j> = 0 for
%            non-consecutive i < j, c) <a_i,a'_{i+1}> in k^*; pp. 103-105 the
%            proof — a triangular (n+2)x(n+2) matrix with diagonal
%            c_{-1}, ..., c_n for necessity; for sufficiency a rank-2 direct
%            factor F, a line D, the line Delta = Ker(F -> L), and s_0 a
%            section of F avoiding D and Delta in every fibre, « ok ».
%   107-112  « Classes de Stiefel-Whitney » of a quadratic form f_V
%            (discriminant (beta_0^n ... beta_{n-1}) rho): the orthogonal
%            projection s~_0, the coefficients xi_i = L_{n-i}(...) = rho_{i+1,n-1},
%            f_V(s~_0) = rho~ sigma~ / rho^2, the recursion
%            w_*(f_{V_{n+1}}) = [1 + d(rho~_{0,n-1} rho~_{1,n-1})] w_*(f_{V_n}),
%            the product formula w_*(f_V) = prod (1 + eta_i + eta_{i+1}) with
%            eta_i = d(rho~_{i,n-1}) in H^1(-, Z/2) (NB w_{n+1} = 0),
%            w_*(f_E); the case n = 2 (w_1, w_2), c(f_V) = w_2 + w_1^2, the
%            Brauer class Br(f_V) = (eps + eta_0 + eta_1 + eta_2)^2 + (sum of
%            eta_i eta_j), the universal base S_0 = Spec Z[beta_0,beta_1]_{...};
%            p. 112 a problem (Stiefel-Whitney classes of the generic smooth
%            quadratic form over Z[(c_ij)]_{2 delta'}) and a NB that w_2, w_3,
%            ... cannot be expressed multiplicatively from d(-1), d(2),
%            d(delta').
%
% Arithmetic checked: the determinant on p. 104, f_V(s~_0) and f(u) =
% beta(4 - beta) (p. 108, correct with phi the polar form of f), the
% identities 4 - beta_0 - beta_1 = -(alpha_0 + alpha_1) etc. with
% beta_i = 2 + alpha_i (p. 110), w_1 and w_2 for n = 2, c(f_V) and Br(f_V)
% mod 2 (pp. 110-111): all agree. One inconsistency is flagged in a \note,
% not corrected: p. 107, the general xi_i at i = 0 gives L_n(beta_1, ...,
% beta_{n-1}) = rho_{1,n-1}, while the line for xi_0 has beta_n and rho_{1,n}.
%
% Notation, fixed once. The gothic letter of pp. 101 and 103 is read
% \mathfrak{h}_i (with H_j for the hyperplanes); E^\vee the dual; his
% barred condition labels (A-bar), (B-bar'), (17-bar) etc. as \bar{} or
% \overline{}; \partial for the Bockstein H^0(-, G_m) -> H^1(-, mu_2) he
% names on p. 107; \mathcal{O}_S for his script O. Long cancelled passages
% (pp. 101-104, under diagonal strokes) are given as \struck{} where they
% read, with one \note each.
%
% Readings to check first: « (24) » at the head of p. 101 and the phrase
% before « k a_i (+) k a_{i+1} = h_i + h_{i+1} »; the rank « n+2 » and the
% index ranges of (9') and (12') on p. 102 (and « base de E » for the a'_i);
% the cancelled lines of pp. 103-104; « (18) » versus « (17) » on pp. 104-105;
% the factor after the parenthesis of delta(f_V) on p. 107; the subscript
% V_n on p. 108; « syst. orthogonal » on p. 110; « la forme x^2+y^2+z^2 » on
% p. 111; « Q_2 » and « H^0(S_0, G_m)_2 » on p. 112.

\documentclass[11pt,a4paper]{article}
\input{../preamble/grothendieck}

\folder{72}
\batch{6}
\pages{101}{112}
\dating{[à partir de 1978]}
\watermark{Édition de démonstration}
\foldertitle{Polyèdres réguliers et hyperpolyèdres réguliers (Calculs en dimension quelconque) : notes manuscrites (s.d.).}

\begin{document}

\page{101}
Pour \uncertain{(24)}, on voit que
\[
z_i z_{i+1} - \rho\, z_{i+1} z_i = \underbrace{\langle a_{i+1}, a'_i\rangle}_{\in k^*} a'_{i+1}\otimes a_i - \rho \underbrace{\langle a_i, a'_{i+1}\rangle}_{\in k^*} a'_i \otimes a_{i+1} \qquad (\rho \in k^*)
\]
\note{les deux « $\in k^*$ » sont écrits sous les crochets, reliés à eux par une flèche montante ; « $\rho \in k^*$ » est écrit sous le membre de gauche}
\struck{qui doit} est $\neq 0$ \struck{\ill{} \ill{}} la famille \struck{donc $\lbrace a'_{i+1}\otimes a_i,\ a'_i\otimes a_{i+1}\rbrace$ est \emph{libre} dans $\operatorname{End}(E)$ — et l'image par $\operatorname{End}(E) \to E$, $u \mapsto u(x)$ formée de $\lbrace a'_{i+1}(x)\, a_i,\ a'_i(x)\, a_{i+1}\rbrace$, qui est libre si} l'image est formée des
\note{le passage de « donc » à « est libre si » est encadré à gauche et annulé par quatre longs traits obliques ; la lecture en est incertaine}
\[
\langle a_{i+1}, a'_i\rangle\, a'_{i+1}(x)\, a_i - \rho \langle a_i, a'_{i+1}\rangle\, a'_i(x)\, a_{i+1},
\]
\ill{} que \ill{} \uncertain{$k a_i \oplus k a_{i+1}$} $= \mathfrak{h}_i + \mathfrak{h}_{i+1}$, \struck{ce} qui \struck{prouve} signifie aussi que
\note{la lettre gothique est lue $\mathfrak{h}$ ; elle revient à la page 102}
\[
(a'_i, a'_{i+1}) : E \longrightarrow k^2 = k e_1 \oplus k e_2
\]
est surjectif, en effet \struck{l'image de} sa restriction à \uncertain{$k a_{i+1} \oplus k a_i$} a la matrice \struck{$a_{i+1} \mapsto (\langle a_{i+1}, a'_i\rangle, 0)$, donc l'image de $E$ contient $k e_1$, en fait l'image}
\note{sous $\langle a_{i+1}, a'_i\rangle$ dans la ligne biffée, « $\in k^*$ » avec une flèche montante}
\[
\begin{pmatrix} \langle a_{i+1}, a'_i\rangle & * \\ 0 & \langle a_i, a'_{i+1}\rangle \end{pmatrix}
\]
de déterminant $\langle a_{i+1}, a'_i\rangle \langle a_i, a'_{i+1}\rangle \in k^\times$, en vertu de (25), lui-même conséquence immédiate de (13) $\lambda_i \in k^\times$, compte tenu de (4). Enfin (25$'$) est une conséquence immédiate de (25),

cqfd.

\page{102}
\emph{Remarque.} Si $E$ est loc. libre de rg \uncertain{$n+2$}, alors on trouve des conditions équivalentes $(\bar A)$, \struck{$(B)$} et $(\bar B)$, en renforçant (A), (B), remplaçant (9) de (A), et (1) de (B), respectivement par
\begin{itemize}
\item[(9$'$)] \uncertain{$(a_i)_{0 \leq i \leq n+1}$} forment une base de $E$
\item[(12$'$)] $(a'_i)_{-1 \leq i \leq n}$ forment une base de $E$.
\end{itemize}
\note{ainsi en (12$'$) ; les $a'_i$ étant des formes linéaires, on attendrait le dual de $E$}
\struck{Quand ces conditions} Elles sont aussi équivalentes à la conjonction de (A), de (B), avec la condition
\[
(26) \qquad E_{n+1} = E .
\]
\struck{Sont alors $\xi'_0 = a'_{-1}$} \struck{Si ces conditions sont satisfaites, complétons}

\struck{Soit $a'_{-1}$ [$\xi'_0 = a'_{-1}$] une forme lin. sur $E$ telle que (27) $\operatorname{Ker} a'_{-1} = k(a_0, \dots, a_n)$ elle est donc déterminée, au facteur scalaire $\in k^*$ près, par la donnée des $z_i$ $(0 \leq i \leq n)$, de \ill{}}
\note{ce passage est encadré à gauche et annulé par des traits obliques croisés ; le numéro « (27) » est d'une lecture incertaine (on pourrait lire « (47) »)}

Idem pour cor. 1 [conditions $(\bar A')$, $(\bar B')$, équivalentes à (A$'$), (B$'$) + (26)]

\emph{Proposition.} Soient $E$ un $k$-module loc. libre de rg $n+2$, $z_i$ $(0 \leq i \leq n)$ des pseudo-réflexions dans $E$,
\[
z_i = \mathrm{id} + a'_i \otimes a_i \qquad (0 \leq i \leq n),
\]
(27) \struck{satisfaisant} \struck{$z_i z_j = z_j z_i$ dans $\underline{\check P}(E)$.}
(28) Pour qu'il existe \struck{un} (localement) $s_0 \in E$ tel que les conditions $(\bar A)$, $(\bar B)$ du \ill{}) (remarque ci-dessus), \struck{soient} satisfaites (conditions de commutation et de « non-dég »), il f. et il s.
\note{« commutation et » est écrit au-dessus de la ligne et appelé par un trait}

\page{103}
que les conditions suivantes soient satisfaites.
\begin{itemize}
\item[a)] $(a_i)_{0 \leq i \leq n}$ est \struck{libre} strictement libre dans $E$ (i.e. engendre un sous-Module facteur direct de rg $n+1$)
\item[a$'$)] $(a'_i)$ est strictement libre dans $E^{\vee}$
\item[b)] $\langle a_i, a'_j\rangle = 0$ \quad $0 \leq i < j \leq n$, $j \neq i+1$ (ou encore $\mathfrak{h}_i \subset H_j$ pour ces $i, j$)
\item[c)] \struck{\ill{}} $\langle a_i, a'_{i+1}\rangle \in k^*$ \quad $(0 \leq i \leq n-1)$ ($\Longleftrightarrow$ $\mathfrak{h}_i \not\subset H_{i+1}$ en chaque fibre, $0 \leq i \leq n-1$)
\end{itemize}
[Les conditions impliquent \struck{donc} (28) $\langle a_i, a'_j\rangle = 0$ pour $0 \leq i, j \leq n$, $i, j$ non consécutifs).

\emph{Dém.} La nécessité de a) b) c) est claire, \struck{prouvons} \struck{celle de a$'$), prouvons \ill{}} \struck{prouvons \ill{} conditions du th. 1 \ill{} propriété de stabilité des \ill{} par dualité. Ces conditions $E_n \subset E_{n+1} = E$ \ill{} hyperplan, et soit $s'_0$ une}
\note{« c) » est ajouté au-dessus de la ligne ; la suite, jusqu'à « soit $s'_0$ une », est biffée et, pour l'essentiel, annulée par des traits obliques ; la lecture en est très incertaine}

On suppose donc que $s_0$ satisfait les conditions dites, \struck{on pose $a_{-1} = s_0$, et on définit}
\[
a_{-1} = s_0 \in E, \qquad a'_{n+1} \in E^{\vee}
\]
défini par
\[
\operatorname{Ker} a'_{n+1} = k(a_{-1}, a_0, \dots, a_{n-1})
\]
\struck{d'où}
\[
\begin{cases} \langle a_i, a'_{n+1}\rangle = 0 & -1 \leq i \leq n-1 \\ \langle a_n, a'_{n+1}\rangle \in k^* \end{cases}
\]
\note{devant la seconde ligne, un début biffé « $c_n =$ » surmonté d'un mot illisible}
Considérons
\[
\begin{aligned} &(a_{-1}, a_0, \dots, a_n) \in E^{n+2} \\ &(a'_0, a'_1, \dots, a'_{n+1}) \in E^{\vee\, n+2} \end{aligned}
\]
et la matrice qu'ils forment, i.e. une matrice triangulaire, avec sur la diagonale

\page{104}
\[
c_{-1} = \langle a_{-1}, a'_0\rangle, \quad c_0 = \langle a_0, a'_1\rangle, \quad \dots, \quad \& \ c_n = \langle a_n, a'_{n+1}\rangle
\]
donc de déterminant égal à
\[
\Delta = c_{-1} c_0 \cdots c_n = \underbrace{\lambda_n}_{\in k^\times} \underbrace{c_n}_{\in k^\times} \in k^\times ,
\]
\note{devant $\Delta$, un « et » biffé ; les deux « $\in k^\times$ » sous $\lambda_n$ et $c_n$ sont reliés à eux par une flèche montante}
ce qui prouve que \struck{les} $(a'_0, \dots, a'_n, a'_{n+1})$ forment une base de $E^{\vee}$, d'où (a$'$).

Inversement, supposons a) a$'$) b) c) satisfaits, cherchons $s \in E$ satisfaisant, avec $z_0, \dots, z_n$, la condition $(\bar B')$(17) pour $i = -1$, \struck{ce qui} équivaut aux conditions (b) de $(\bar B')$. \struck{En} \struck{Tout d'abord}
\note{suit un passage biffé et annulé par des traits en zigzag, dont on ne lit que des fragments : « (28) \ill{} \ill{}, \ill{} Rappelons la situation \ill{} pour \ill{} (17) \ill{} » ; puis, annulé par des traits obliques croisés :}
\struck{La condition (7) est déjà satisfaite par hypothèse, La condition (8) signifie $\langle s_0, a'_i\rangle a_i = 0$ pour $1 \leq i \leq n$, et comme $a_i$ est libre, ceci équivaut à (29) $\langle s_0, a'_i\rangle = 0$ \quad $1 \leq i \leq n$}

— En vertu de a$'$), l'ens. des $s_0$ satisfaisant \uncertain{(18)} pour $i = -1$ (pour $s_0 = a_{-1}$) \struck{ou conditions} est un sous-module facteur direct, loc. libre de rang $(n+2) - n = 2$, \struck{Soit} soit $F$. Pour tenir compte de $(\bar B')$(19) pour $i = 0$, (Soit $D \subset F$ la droite, définie par les conditions $\langle s_0, a'_i\rangle = 0$ \quad $0 \leq i \leq n$ dans $E$,) \struck{donc} on a la seule condition $\langle s_0, a'_0\rangle \in k^\times$ dans $F$. La condition (19) $\lambda_n \in k^*$ signifie compte tenu de (6) et (4), que $\langle s_0, a'_0\rangle \in k^*$, \struck{$\langle s_0, a'_0\rangle \in$} i.e. $s_0 \notin D$ sur chaque fibre.
\note{« Pour tenir compte de $(\bar B')$(19) pour $i = 0$, » est écrit au-dessus de la ligne ; le numéro (18) pourrait être un (17) surchargé}

\page{105}
Reste à exprimer $(\bar B')$\uncertain{(18)}, pour ceci considérons \struck{la condition formes lin. $a'_{-1}$ qui s'annulent sur les} \struck{le module} $L = E / k(a_0, \dots, a_n)$, et l'homom. \struck{$a_i$} induits par $E \to L$, $F \to L$, \struck{on sait} je dis que c'est un épi, i.e. que c'est un épi sur chaque fibre \struck{point}. Comme on est ramené au cas d'un corps de base, il suffit de prouver que dans ce cas \struck{$F$} on n'a pas $F \subset k(a_0, \dots, a_n)$, \struck{ou encore que} \struck{Or} Or $F \cap k(a_0, \dots, a_{n-1}) = (0)$, car le système des $n$ formes linéaires $a'_1, \dots, a'_n$ sur $k(a_0, \dots, a_{n-1})$ a une matrice triangulaire, de déterminant le produit des $c_i = \langle a_i, a'_{i+1}\rangle$ $(0 \leq i \leq n-1)$, donc $\in k^*$. Comme \uncertain{$k(a_0, \dots, a_{n-1})$} est de codim 1 dans \uncertain{$k(a_0, \dots, a_n)$}, \struck{\ill{}} et $\dim F = 2$, cela implique qu'on ne peut avoir $F \subset k(a_0, \dots, a_n)$.
\note{« $(\bar B')$ » est écrit au-dessus de « (18) », qu'un trait oblique barre peut-être ; dans la phrase « Comme … de codim 1 dans … », les derniers indices des deux modules sont mal formés}

Soit \struck{$D'$} $\Delta = \operatorname{Ker}(F \to L)$, qui est une droite de $F$. La condition $(\overline{18})$ sur $s_0$ signifie que $s_0 \notin \Delta$ en chaque fibre. Donc la conjonction de $(\overline{17})$ $(\overline{18})$ $(\overline{19})$ signifie que $s_0$ est une section de $F$ (fibré vect. de rg 2) qui n'est dans $D$ ni $\Delta$ en aucune fibre — il en existe toujours localement, ok.

\section*{Classes de Stiefel-Whitney et classe de Brauer de la forme invariant}

\note{inscrit à l'encre en haut à droite du feuillet de garde (page 106), de papier chamois, qui ne porte rien d'autre ; « invariant » est écrit ainsi. Le feuillet ne reçoit pas de numéro de page}

\page{107}
\marginal{* Mais il faut supposer \uncertain{$f \mid \mathcal{O}_S[a_1, \dots, a_n]$} lisse, donc $\sigma = \beta_n$ inversible}
\note{cette note est écrite en haut à droite de la page ; l'astérisque renvoie à celui qui suit « $\mathcal{O}_S[a_1, \dots, a_n]$ » cinq lignes plus bas}

\emph{Classes de Stiefel-Whitney.}

Supposons $f_V$ lisse \struck{et} 2 inv., i.e. \struck{2}
\[
\delta(f_V) = (\beta_0^{\,n} \beta_1^{\,n-1} \cdots \beta_{n-1})\, \rho \quad \text{inv. et } 2 \text{ inv.}
\]
\note{devant la parenthèse, un chiffre biffé (un 2 ?) ; le facteur final, lu $\rho$, est d'une lecture incertaine}
(Calculons $w_i(f_V)$. Pour ceci, on va procéder par récurrence en calculant l'orthogonal de $\mathcal{O}_S[a_1, \dots, a_n]$\,*, Comme
\[
\Psi_V(a_i) = -(\beta_0 \cdots \beta_{i-1})\, a'_i ,
\]
c'est aussi l'orthogonal de $a'_1, \dots, a'_n$, donc la projection orthogonale \struck{de} $s_0 - \frac{1}{\rho}\pi$ de $s_0 = a_{-1}$ sur $V$, dans $E = V \oplus \mathcal{O}_S \pi$, soit $p = \mathrm{id} - \frac{1}{\rho}\, \pi' \otimes \pi$ la dite projection [notée $\tilde s_0$]
\note{« orthogonale » et « $s_0 - \frac{1}{\rho}\pi$ » sont écrits au-dessus de la ligne ; le $\mathcal{O}$ de $\mathcal{O}_S \pi$ est souligné}
\[
[\ \tilde s_0 = -\frac{1}{\rho}\bigl(\xi_0 a_0 + \xi_1 a_1 + \cdots + \xi_n a_n\bigr)
\]
\[
\begin{cases}
\xi_n = 1 \\
\xi_{n-1} = 2 \\
\xi_i = L_{n-i}(\beta_{i+1}, \dots, \beta_{n-1}) = \rho_{i+1,\,n-1} \\
\xi_0 = L_n(\beta_1, \dots, \beta_n) = \rho_{1,n} = \sigma
\end{cases}
]
\]
\note{le $\xi_0$ de la dernière ligne est écrit par-dessus une autre lettre. La formule générale donnerait pour $i = 0$ : $L_n(\beta_1, \dots, \beta_{n-1}) = \rho_{1,n-1}$, alors que la dernière ligne porte $\beta_n$ et $\rho_{1,n}$ ; l'écart est sur la page}
\[
f_V(\tilde s_0) = f_E\Bigl(s_0 - \frac{1}{\rho}\pi\Bigr) = \underbrace{f_E(s_0)}_{0} + \frac{1}{\rho^2} \underbrace{f_E(\pi)}_{-\tilde\rho\tilde\sigma} - \frac{1}{\rho}\, \underbrace{\varphi(s_0, \pi)}_{\langle s_0,\, \psi(\pi)\rangle \,=\, \langle s_0,\, -\sigma\pi'\rangle \,=\, -\sigma}
\]
\note{le signe devant $\frac{1}{\rho}\varphi(s_0, \pi)$ est surchargé ; la ligne suivante le lit comme un moins. Sous $\varphi(s_0, \pi)$ il écrit $\langle s_0, \psi(\pi)\rangle$, avec $\psi(\pi)$ souligné et $-\sigma\pi'$ en dessous, puis $-\sigma$ sous une accolade ; la mise en ligne est nôtre}
\[
= -\frac{\tilde\rho\tilde\sigma}{\rho^2} + \frac{\sigma}{\rho} = \frac{1}{\rho^2}\Bigl[\underbrace{\rho\sigma}_{2\tilde\rho\tilde\sigma} - \tilde\rho\tilde\sigma\Bigr] = \frac{\tilde\rho\tilde\sigma}{\rho^2}
\]
\[
\partial\bigl(f_V(\tilde s_0)\bigr) = \partial(\tilde\rho\tilde\sigma) = \partial(\tilde\rho_{0,n-1}\, \tilde\rho_{1,n-1})
\]
\marginal{sous le premier $\partial$ : Bockstein $H^0(\ ,\, \mathbb{G}_m) \to H^1(\ ,\, \mu_2)$}

\page{108}
Donc on trouve
\[
w_*\bigl(f_{V_{n+1}}(\beta_0, \dots, \beta_{n-1})\bigr) = \bigl[1 + \underbrace{\partial(\tilde\rho_{0,n-1}\, \tilde\rho_{1,n-1})}_{\in H^1(-,\, \mu_2)}\bigr]\, w_*\bigl(f_{V_n}(\beta_1, \dots, \beta_{n-1})\bigr)
\]
\note{l'indice du second $V$ est mal formé (lu $V_n$) ; dans $H^1(-, \mu_2)$, le coefficient est écrit par-dessus un autre, biffé}
On trouve ainsi de proche en proche, pourvu que les calculs aient un sens i.e. que \struck{tous} les $\tilde\rho_{i,n-1}$ soient inversibles $(0 \leq i \leq n-1)$
\[
\begin{aligned}
w_*\bigl(f_{V_{n+1}}(\beta_0, \dots, \beta_{n-1})\bigr) = {}& \bigl[1 + \partial(\tilde\rho_{0,n-1}\, \tilde\rho_{1,n-1})\bigr] \bigl[1 + \partial(\tilde\rho_{1,n-1}\, \tilde\rho_{2,n-1})\bigr] \\
& \cdots \bigl[1 + \partial(\tilde\rho_{n-2,n-1}\, \tilde\rho_{n-1,n-1})\bigr]\, w_*\bigl(f_{V_2}(\beta_{n-1})\bigr)
\end{aligned}
\]
\note{dans le dernier crochet, les indices sont surchargés et récrits en dessous : $n-2, n-1$ et $n-1, n-1$}
De même, on trouve,
\[
w_*\bigl(f_{V_2}(\beta_{n-1})\bigr) = \bigl(1 + \partial[\underbrace{\tilde\rho_{n-1,n-1}}_{4 - \beta_{n-1}}\, \underbrace{\tilde\rho_{n,n-1}}_{1}]\bigr) \underbrace{\bigl(1 + \partial(\underbrace{\tilde\rho_{n,n-1}}_{1}\, \underbrace{\tilde\rho_{n+1,n}}_{1})\bigr)}_{1}
\]
\[
\bigl(1 + \partial(\beta(4 - \beta))\bigr) = \bigl[1 + \partial\bigl((2 + \alpha)(2 - \alpha)\bigr)\bigr]
\]
\note{un long trait oblique traverse le membre de droite de la première ligne, qui semble remplacé par la seconde ; sous $\tilde\rho_{n,n-1}$ du premier facteur, un signe de répétition ; entre $(2+\alpha)$ et $(2-\alpha)$, un signe surchargé}
\struck{En} En effet, faisons le calcul pour $f_V(\beta)$ \ill{} $a_0, a_1$ \ill{} \ill{}, si on veut
\[
\begin{cases}
f(a_0) = 1 \qquad \partial f(a_0) = 0 \\
a_0^{\perp} = \lbrace x a_0 + y a_1 \mid 2x - \beta y = 0 \ \text{ i.e. } \ x = \frac{\beta}{2} y \rbrace
\end{cases}
\]
\note{devant « $x = \frac{\beta}{2} y$ », un « $y =$ » biffé}
engendré par $\beta a_0 + 2 a_1 = u$
\[
\begin{aligned}
f(u) &= \beta^2 \underbrace{f(a_0)}_{1} + 4 \underbrace{f(a_1)}_{\beta} + 2\beta \underbrace{\varphi(a_0, a_1)}_{-\beta} \\
&= 4\beta - \beta^2 = \beta(4 - \beta)
\end{aligned}
\]
\note{devant $\beta^2 f(a_0)$, un début biffé « $\beta^2 + 4\beta$ ». Le calcul est juste pour $\varphi$ forme polaire de $f$}
d'où \struck{$w_*$} la formule annoncée.

\page{109}
On trouve, en résumé,
\[
w_*\bigl(f_V(\beta_0, \dots, \beta_{n-1})\bigr) = \Bigl(\prod_{0 \leq i \leq n-2} \bigl(1 + \partial(\tilde\rho_{i,n-1}\, \tilde\rho_{i+1,n-1})\bigr)\Bigr) \Bigl(1 + \partial\bigl(\underbrace{\beta_{n-1}(4 - \beta_{n-1})}_{(2 + \alpha_{n-1})(2 - \alpha_{n-1})}\bigr)\Bigr)
\]
formule valable si tous les $\tilde\rho_{i,n-1}$ $(0 \leq i \leq n-1)$ sont inversibles, ainsi que les $\beta_i$ $(0 \leq i \leq n-1)$.

Bien sûr, on aurait pu conduire le calcul en sens inverse, en \ill{} \ill{} partant de $(a_0, \dots, a_{n-1})$ etc, \struck{on} trouverait-on d'autres expressions remarquables ?

Calculons $w_*(f_E)$, \struck{si on suppose} on sait que $V^{\perp} = \mathcal{O}_S \pi$ et $f(\pi) = -\tilde\rho\tilde\sigma$ et on trouve donc un facteur supplémentaire
\[
1 + \partial(\underbrace{-\tilde\rho\tilde\sigma}_{-\tilde\rho_{0,n-1}\, \tilde\rho_{1,n-1}}) .
\]
Posons donc
\[
\begin{aligned}
\partial(-1) &= \varepsilon \\
\partial(\tilde\rho_{i,n-1}) &= \eta_i \qquad 0 \leq i \leq n-1 \\
\partial(\beta_{n-1}) &= \eta_n
\end{aligned}
\qquad \text{éléments de } H^1(-, \mathbb{Z}/2)
\]
\note{un tilde se trouve entre $\tilde\rho_{i,n-1}$ et $\beta_{n-1}$ ; il peut appartenir à l'une ou l'autre lettre}
On trouve
\[
w_*(f_V) = \prod_{0 \leq i \leq n-1} (1 + \eta_i + \eta_{i+1})
\]
\marginal{NB $w_{n+1}(f_V) = 0$ !}
\[
w_*(f_E) = (1 + \varepsilon + \eta_0 + \eta_1) \prod_{0 \leq i \leq n-1} (1 + \eta_i + \eta_{i+1})
\]
\note{dans la première formule, un signe biffé avant le « $=$ »}

\page{110}
Si $n = 2$, ceci donne
\[
\begin{cases}
\varepsilon = \partial(-1) \\
\eta_0 = \partial(\tilde\rho) = \partial(\underbrace{4 - \beta_0 - \beta_1}_{-(\alpha_0 + \alpha_1)}) = \partial\bigl[-(\alpha_0 + \alpha_1)\bigr] \\
\eta_1 = \partial\tilde\sigma = \partial(4 - \beta_1) = \partial(2 - \alpha_1) \\
\eta_2 = \partial\beta_1 = \partial(2 + \alpha_1)
\end{cases}
\]
\[
w_* f_V = (1 + \eta_0 + \eta_1)(1 + \eta_1 + \eta_2)
\]
avec hyp. parasite que $\tilde\sigma = 2 - \alpha$, \struck{\ill{}} soit inversible. Un calcul plus satisfaisant, sans cette hypothèse, s'obtient en prenant l'orthogonal de $\mathcal{O}_S(a_0, a_2)$ engendré par $u = \beta_0 a_0 + 2 a_1 + a_2$, on trouve $f(u) = -\beta_0(\alpha_0 + \alpha_1)$, et $f(a_0) = 1$, $f(a_2) = \beta_0\beta_1$ \uncertain{syst. orthogonal}. \ill{} pose donc
\note{« soit » est écrit au-dessus d'un mot biffé ; le signe entre $\tilde\sigma$ et $2 - \alpha$ est surchargé}
\[
\begin{aligned}
\varepsilon &= \partial(-1) \\
\eta_0 &= \partial(\beta_0) \\
\eta_1 &= \partial(\beta_1) \\
\eta_2 &= \partial\bigl[-(\alpha_0 + \alpha_1)\bigr]
\end{aligned}
\]
et on trouve
\[
\begin{aligned}
w_*(f_V) &= (1 + \eta_0 + \eta_2)(1 + \eta_0 + \eta_1) \quad \text{i.e.} \\
w_1(f_V) &= \eta_1 + \eta_2 \\
w_2(f_V) &= (\eta_0 + \eta_1)(\eta_0 + \eta_2) = \eta_0^2 + (\eta_0\eta_1 + \eta_1\eta_2 + \eta_2\eta_0)
\end{aligned}
\]
\note{dans les deux dernières lignes, un début biffé et illisible après le « $=$ »}

\page{111}
On a
\[
c(f_V) \overset{\text{dfn}}{=} w_2(f_V) + w_1(f_V)^2 = (\eta_0 + \eta_1 + \eta_2)^2 + (\eta_0\eta_1 + \eta_1\eta_2 + \eta_2\eta_0)
\]
(La condition n. et s. pour que le fibré en droites projectives associé soit isomorphe à celui défini par \struck{\ill{}} \uncertain{la forme} $x^2 + y^2 + z^2$). Donc
\[
\operatorname{Br}(f_V) \overset{\text{dfn}}{=} c(f_V) - \underbrace{c(X^2 + YZ)}_{\varepsilon^2} = (\varepsilon + \eta_0 + \eta_1 + \eta_2)^2 + (\eta_0\eta_1 + \eta_1\eta_2 + \eta_2\eta_0)
\]
\note{dans $c(X^2 + YZ)$, un mot biffé devant $X^2$}
L'annulation de cette expression est n. et s. pour que le fibré en droites projectives \struck{associé} défini par $f_V$ soit associé à un fibré vect. de rang 2, en fait, c'est la classe de Brauer de ce fibré … Le calcul universel se fait sur
\[
S_0 = \operatorname{Spec} \mathbb{Z}[\beta_0, \beta_1]_{2\beta_0\beta_1(4 - \beta_0 - \beta_1)}
\]
On notera que dans $H^1(S_0, \mathbb{Z}/2)$, les \struck{sections} éléments $\varepsilon, \eta_0, \eta_1, \eta_2$ forment une base. Il est plausible que les 10 expressions
\[
\varepsilon^2,\ \eta_0^2,\ \eta_1^2,\ \eta_2^2,\ \varepsilon\eta_0,\ \varepsilon\eta_1,\ \varepsilon\eta_2,\ \eta_0\eta_1,\ \eta_1\eta_2,\ \eta_2\eta_0
\]
sont
\note{dans $H^1(S_0, \mathbb{Z}/2)$, l'exposant est écrit par-dessus un autre chiffre}

\page{112}
? lin. indép. dans $H^2(S_0, \mathbb{Z}/2)$ …
\note{le point d'interrogation est dans la marge de gauche, en tête de page}

\emph{Pb.} Considérons une forme quadratique \add{lisse} à coeff. indéterminés, à coefficients dans
\[
\mathbb{Z}\bigl[(c_{ij})_{1 \leq i \leq j \leq n}\bigr]_{2\delta'((c_{ij}))}
\]
« Calculer » ses classes de Stiefel-Whitney totales ?

\emph{NB} Le $H^1(S_0, \mathbb{Z}/2)$ ($\simeq$ \uncertain{$H^0(S_0, \mathbb{G}_m)_2$}) est engendré par
\[
\varepsilon = \partial(-1), \quad \underbrace{\partial(2)}_{\xi}, \quad \underbrace{\partial(\delta'(c_{ij}))}_{\eta \,=\, w_1(f)},
\]
\note{il met « $w_1(f)$ » entre guillemets ; « $\simeq H^0(S_0, \mathbb{G}_m)_2$ » est écrit au-dessus de « est engendré par » ; devant $\partial(2)$, un signe biffé}
dont les 2 premiers deviennent nuls sur une extension biquadratique \uncertain{$\mathbb{Q}_2$}$[\sqrt{-1}, \sqrt{2}]$ — a fortiori sur $\mathbb{C}$, le deuxième donne lieu à $w_1$. Il est donc clair que multiplicativement à partir de ces éléments, on ne peut exprimer aucun des $w_2, w_3, \dots$ etc.
\note{la lettre du corps de base est surchargée (un $\mathbb{Z}$ barré ?) ; « donc » et « aucun des » sont écrits au-dessus de la ligne}

\end{document}
